\section{Исследование и построение решения задачи}
\label{sec:Chapter3} \index{Chapter3}

В разделе из общей содержательной задачи постановки
последовательно выделяются подзадачи, и для каждой из них строится
решение. При этом на каждом шаге фиксируются принятые проектные решения и
обоснование выбора.

\subsection{Декомпозиция задачи}

Поставленная задача --- построить отображение
$\pi: \{M_k\} \mapsto \{M'_j\}$ --- разбита на следующие
подзадачи:
\begin{enumerate}
    \item \textbf{Pipeline формирования BLAS для динамических
      ассетов в DagorEngine.} Без точного знания цепочки
      \texttt{mesh}, \texttt{geometry-array}, \texttt{BLAS}
      невозможно определить, какое представление меша менять
      и в какой момент времени.
    \item \textbf{Алгоритм объединения и переразбиения
      треугольников модели.} Должен сохранять множество
      треугольников, перегруппировывать их по пространственной
      близости и удовлетворять аппаратным ограничениям DXR
      (16-битный индекс, лимит числа геометрий на BLAS).
    \item \textbf{Способ валидации гипотезы.} Необходимо
      экспериментально подтвердить, что метод даёт реальное
      ускорение GPU-трассировки в работающем семпле движка.
\end{enumerate}

Подзадачи решаются последовательно в
подразделах ниже.

\subsection{Pipeline формирования BLAS в DagorEngine}
\label{ssec:p1}

Исследование исходного кода DagorEngine~\cite{dagorEngine} показывает,
что при подготовке BLAS для динамической модели выполняется следующая
цепочка преобразований:

\begin{equation}
\begin{split}
&\underbrace{\texttt{DynamicRenderableSceneLodsResource}}_{\text{ресурс модели}} \\
&\quad\;\to\; \underbrace{\texttt{ShaderMesh::RElem[]}}_{\text{меш-уровень}} \\
&\quad\;\to\; \underbrace{\texttt{RaytraceGeometryDescription[]}}_{\text{DXR-вход}} \\
&\quad\;\to\; \underbrace{\mathcal{B}}_{\text{драйвер}} \;\to\; \underbrace{T}_{\text{BLAS}}.
\end{split}
\label{eq:pipeline}
\end{equation}

\noindent
\textbf{(1) \texttt{DynamicRenderableSceneLodsResource}}~---~ресурс
динамической модели в Dagor; содержит несколько уровней детализации, каждый из
которых --- набор \texttt{ShaderMesh}. Для построения BLAS
используется наивысший уровень.

\noindent
\textbf{(2) \texttt{RElem}}~---~атомарная меш-структура графического движка. Содержит
указатель на материал, ссылку на разделяемый вершинный буфер
\texttt{vertexData} и диапазоны \texttt{(sv, numv, si, numf,
baseVertex)}. Один \texttt{RElem}~=~непрерывный поддиапазон
индексов с единой шейдер-материальной атрибуцией.

\noindent
\textbf{(3) \texttt{RaytraceGeometryDescription}}~---~DXR-совместимая
структура.
Массив таких структур и является \emph{основным входом}
вендорского \texttt{build\_acceleration\_structure}. Каждая запись описывает
одну geometry (vertex buffer, index buffer, диапазон, формат,
флаги).

\noindent
\textbf{(4) Билдер $\mathcal{B}$}~---~внутренний код графического
драйвера. В соответствии со спецификацией DXR~\cite{dxrSpec}
обязан сохранять
$(\texttt{GeometryIndex},\,\texttt{PrimitiveIndex})$ как стабильный
идентификатор треугольника. Это формальное ограничение запрещает
драйверу выполнять реорганизацию между geometry: например,
объединять треугольники двух \texttt{Description}-ов в один лист
BVH.

\paragraph{Точка входа BVH-подсистемы.} Цепочка формирования BLAS
обслуживается функциями \texttt{bvh::add\_mesh()} и
\texttt{bvh::add\_instance()} из
\texttt{prog/gameLibs/bvh/bvh.cpp}. При загрузке динамической
модели исследуемый семпл \texttt{testGI} формирует BLAS-ы
на уровне целых ассетов: для каждого \texttt{RElem} строится свой \emph{BLAS}.
Именно на этой стадии, перед вызовом построения, в семпл встраивается предлагаемый метод предобработки.

\paragraph{Аппаратное ограничение 16-битного индекса.} Внутри
\texttt{bvh.cpp} зафиксировано условие
\texttt{G\_ASSERT(mesh.indexFormat == 2)} --- BLAS принимает только
16-битные индексы. Это означает, что одна
\texttt{RaytraceGeometryDescription} может ссылаться максимум на
$2^{16}\!-\!1 = 65\,535$ различных вершин; при 
худшем случае --- на не более чем
$\lfloor 65\,535 / 3 \rfloor = 21\,845$ треугольников. Метод
обязан учитывать это ограничение при разбиении модели.

\paragraph{Зона вмешательства.} Из этой цепочки следует,
что \emph{практически применимая} точка, доступная разработчику без
модификации драйвера, --- стадия формирования
\texttt{RElem[]} и \texttt{RaytraceGeometryDescription[]}.
Эта стадия выполняется в коде \texttt{testGI} непосредственно
перед вызовом \texttt{bvh::add\_mesh()}. Преобразование $\pi$
встраивается как hook на этой границе: вместо подачи в билдер
исходного списка \texttt{RElem} семпл собирает данные всех
\texttt{RElem}-ов модели, перегруппирует их по пространственной
близости и порождает новый набор BLAS-ов с собственными
vertex/index-буферами.

\subsection{Алгоритм глобальной пространственной кластеризации}
\label{ssec:p2}

\subsubsection{Идея}

Гипотеза работы, сформулированная в постановке задачи, связывает проблему со
способом группировки треугольников в ассетах: BLAS строится для всего
ассета целиком, поэтому логическая группа (например, все деревянные
конструкции), разбросанная по объёму модели, получает огромный AABB, а
множество таких ассетов даёт сильно перекрывающиеся инстансы в TLAS.
Поскольку драйвер не может сам разрезать BLAS, единственный совместимый
с DXR способ состоит в том, чтобы на этапе подготовки данных разбить крупный ассет на
множество мелких пространственно-компактных мини-ассетов.

Идея метода:
\begin{enumerate}
    \item \emph{Объединить} треугольники всех \texttt{RElem}-ов
      одной модели в общий пул, забыв временно про
      материальную атрибуцию.
    \item \emph{Глобально отсортировать} пул по пространственной
      близости --- 63-битному Morton-коду центроида треугольника.
    \item \emph{Разрезать} отсортированный пул на $K$ последовательных
      сегментов одинакового размера; каждый сегмент образует
      пространственно-плотный «кусок» модели благодаря
      Z-order-локальности.
    \item \emph{Для каждого сегмента} выделить собственные
      \texttt{vertex} и \texttt{index} буферы, зарегистрировать
      сегмент как отдельный BLAS через
      \texttt{bvh::add\_mesh()}.
    \item \emph{В TLAS} разместить $K$ инстансов новых мини-ассетов.
\end{enumerate}

Наглядно разница между исходным и предлагаемым представлением
показана на схеме ниже: вместо нескольких BLAS с
AABB размером со всю модель метод порождает множество мелких
пространственно-плотных кластеров с почти не перекрывающимися
AABB.

\begin{figure}[h!]
\centering
\includegraphics[width=0.8\textwidth]{diagrams/scheme_idea.png}
\caption{Идея метода. \textbf{Слева:} baseline --- каждая исходная
группа \texttt{RElem} даёт BLAS с AABB во весь дом, инстансы в TLAS
сильно перекрываются ($9{,}15$~мс на худшем, случайном разбиении).
\textbf{Справа:} после Morton-кластеризации дом разбит на $K$ компактных
кластеров с малыми непересекающимися AABB, и время падает до
$1{,}07$~мс ($\mathbf{\times 8{,}6}$).}
\label{fig:idea}
\end{figure}

В этой версии метода материальная атрибуция треугольников теряется
(все треугольники объединяются под унифицированный
ray-tracing-материал), что приемлемо для трассировки теней RTSM,
где материал в hit-shader-е не используется. Это ограничение
обсуждается в заключении.

\subsubsection{Шаги алгоритма}

Пусть исходная модель содержит $R$ \texttt{RElem}-ов с числом
треугольников $|T_1|, \ldots, |T_R|$; обозначим
$N\!=\!\sum_k |T_k|$ --- полное число треугольников модели.

\begin{enumerate}
    \item \textbf{Сбор треугольников.} Для каждого
      \texttt{RElem}-а через \texttt{lockVbufferEx} и
      \texttt{lockIbufferEx} прочитываются позиции вершин
      (формат \texttt{VSDT\_FLOAT3}) и индексы (формат
      \texttt{uint16\_t}). Из них формируется массив
      $\{(\mathbf{p}_0^i, \mathbf{p}_1^i, \mathbf{p}_2^i)\}_{i=1}^{N}$
      позиций треугольников в локальной системе модели.
    \item \textbf{Глобальный AABB и центроиды.} Здесь $t_i$ --- $i$-й
      треугольник модели (с вершинами $\mathbf{p}_0^i, \mathbf{p}_1^i,
      \mathbf{p}_2^i$ из шага 1), а $M$ --- вся модель, то есть
      объединение всех её треугольников $t_1,\ldots,t_N$. Вычисляется
      габаритный бокс всей модели как объединение боксов отдельных
      треугольников
      $\mathrm{AABB}(M)\!=\!\bigcup_i \mathrm{AABB}(t_i)$ и набор
      центроидов (центров масс) треугольников
      $\mathbf{c}_i = (\mathbf{p}_0^i+\mathbf{p}_1^i+\mathbf{p}_2^i)/3$.
      Эти величины --- подготовка к Morton-сортировке (шаг 3):
      $\mathrm{AABB}(M)$ задаёт диапазон нормировки координат в
      $[0,1]^3$, а центроиды $\mathbf{c}_i$ служат ключом сортировки
      треугольников по пространственной близости.
    \item \textbf{Morton-сортировка.} Для каждого центроида
      вычисляется 63-битный Morton-код:
      нормированные координаты $\hat{\mathbf{c}}_i \in [0,1]^3$
      квантуются в 21-битные целые $\mathbf{q}_i \in
      \{0,\ldots,2^{21}-1\}^3$, и биты перемежаются:
      \[
          \mu(\mathbf{c}_i) =
          \sum_{b=0}^{20}\!\Bigl(
              q_{i,x}^{(b)}\!\cdot\!2^{3b}
              + q_{i,y}^{(b)}\!\cdot\!2^{3b+1}
              + q_{i,z}^{(b)}\!\cdot\!2^{3b+2}
          \Bigr).
      \]
      Треугольники сортируются стабильно по убыванию приоритета
      $\mu$.
    \item \textbf{Согласование $K$ с лимитом 16-битного индекса.}
      Пусть пользовательский запрос на число кластеров --
      $K_{\text{req}}$ (параметр
      \texttt{bvhMergeAllRelems} в \texttt{settings.blk}).
      Эффективное число кластеров:
      \[
          K \;=\; \max\!\Bigl(K_{\text{req}},\
              \bigl\lceil N / T_{\max} \bigr\rceil\Bigr),
          \qquad T_{\max} = 21\,845.
      \]
      Условие $K \ge \lceil N/T_{\max} \rceil$ гарантирует, что
      ни один кластер не превысит лимит 16-битного индекса.
    \item \textbf{Разбиение и формирование BLAS-ов.}
      Отсортированный пул делится на $K$ последовательных
      сегментов размера $\sim N/K$. Для каждого сегмента $j$:
      \begin{itemize}
          \item выделяется собственный vertex-буфер $V_j$;
          \item выделяется собственный index-буфер $I_j$;
          \item index-буфер создаётся с флагами
            \texttt{SBCF\_BIND\_SHADER\_RES} и
            \texttt{SBCF\_MISC\_ALLOW\_RAW}; они нужны для
            обращения к буферу как к SRV в compute-шейдере
            \texttt{BVH::IndexProcessor};
          \item сегмент регистрируется через
            \texttt{bvh::add\_mesh()} с уникальным
            \texttt{mesh\_id}, образованным из идентификатора
            ресурса и индекса кластера.
      \end{itemize}
    \item \textbf{Размещение инстансов.} При вызове
      \texttt{bvh::add\_instance()} в TLAS добавляется $K$
      инстансов новых кластеров с общей матрицей трансформации
      экземпляра.
\end{enumerate}

Полный конвейер из шести шагов и его единственный управляющий
параметр сведены на схеме конвейера ниже.

\begin{figure}[h!]
\centering
\includegraphics[width=0.85\textwidth]{diagrams/scheme_pipeline.png}
\caption{Конвейер метода: от сбора \texttt{RElem}-ов модели через
глобальную Morton-сортировку и разбиение на $K$ кластеров до регистрации
их в BVH. Управление --- одна строка \texttt{bvhMergeAllRelems} в
\texttt{settings.blk}; эффективное число кластеров согласуется с
аппаратным лимитом 16-битного индекса.}
\label{fig:pipeline}
\end{figure}

\subsubsection{Сложность}

Сбор треугольников и вычисление центроидов --- $O(N)$ при пиковой
памяти $O(N)$. Morton-сортировка стабильна и в реализации
основана на \texttt{eastl::sort}; в худшем случае
$O(N\log N)$. Формирование $K$ vertex/index-буферов и
регистрация --- $O(N)$. Итоговая асимптотика метода --
$O(N\log N)$, пиковая память $O(N)$. Алгоритм работает один раз
при инициализации BVH-контекста и не влияет на per-frame
стоимость рендера.

\subsubsection{Корректность}

\begin{enumerate}
    \item Множество треугольников не меняется. Каждый треугольник
      исходного меша попадает ровно в один кластер. Это
      обеспечивает покрытие сцены в трассировке.
    \item Позиции вершин копируются побитово (с приведением
      исходного формата к \texttt{Point3} при необходимости);
      визуальная эквивалентность теневой картинки следует из
      инвариантности набора треугольников.
    \item Автоматическое
      увеличение $K$ до $\lceil N/T_{\max}\rceil$ гарантируют
      работоспособность метода в рамках аппаратных ограничений
      DXR Tier 1.1.
\end{enumerate}

\subsection{Способ валидации}
\label{ssec:p3}

Улучшение качества BVH должно отражаться на реальной
производительности, иначе ценность метода сомнительна.
Поэтому валидация проводится в семпле \texttt{testGI} DagorEngine, который
рассчитан на нагрузочные эксперименты с рейтрейсингом и обладает
встроенным GPU-профайлером с per-pass таймингами.

\paragraph{Тестовая модель.} За основу взята общедоступная модель
жилого дома
\texttt{House, Ko Bulon Don Urak Lawoi Village, Thailand}. Для изоляции исследуемой проблемы и
проведения контролируемых тестов из неё подготовлены два варианта
разбиения на меш-ассеты:
\begin{itemize}
    \item \textbf{«Плохой» вариант} (случайный split): треугольники
      модели разбиты на группы \emph{случайным образом}, без учёта
      материала и положения в пространстве. В результате каждая группа
      равномерно «размазана» по всему объёму дома, и её \texttt{RElem}
      получает AABB размером со всю модель --- синтетический худший
      случай с максимальными перекрытиями в TLAS.
    \item \textbf{Реалистичный вариант} (материальный split):
      треугольники сгруппированы \emph{по материалам} (металлические
      детали дверей и сантехники в одной группе, дерево в другой,
      стекло в третьей, и т.\,д.), как это обычно делает пайплайн
      игрового движка. Материалы частично локализованы по частям
      здания, поэтому перекрытия AABB слабее, чем при случайном
      разбиении; этот вариант служит честной точкой сравнения,
      соответствующей реальным production-моделям и гипотезе работы.
\end{itemize}
Оба варианта содержат одно и то же множество треугольников и дают
попиксельно одинаковую картинку; при этом различается лишь группировка
треугольников по \texttt{RElem}-ам.

\paragraph{Тестовая сцена.} В сцену \texttt{testGI} размещены
$N\!=\!16$ инстансов модели \texttt{House, Ko Bulon Don Urak Lawoi Village, Thailand}, чтобы получить
заметную RT-нагрузку. Активна подсистема ray-traced shadow maps
для одного направленного источника света (солнце). Итоговый вид
сцены приведён на рисунке ниже.

\begin{figure}[h!]
\centering
\includegraphics[width=0.78\textwidth]{scene.png}
\caption{Тестовая сцена в \texttt{testGI}: сетка $4\!\times\!4$ из
$16$ инстансов модели жилого дома с трассируемыми тенями от солнца. Именно на этой сцене
измеряется время прохода \texttt{rtsm::do\_trace}.}
\label{fig:scene_render}
\end{figure}

\paragraph{Замеряемая величина.} Медианное GPU-время прохода
\texttt{rtsm::do\_trace} из встроенного GPU-профайлера движка.
Замер выполняется в установившемся режиме (тёплый кэш
RTSM-денойзера, неподвижная камера, фиксированный кадровый
буфер).

\paragraph{Сравниваемые конфигурации.}
\begin{itemize}
    \item \textbf{Baseline.} Все опции оптимизации отключены
      (\texttt{bvhMergeAllRelems:i = 0}).
    \item \textbf{Глобальная пространственная кластеризация.} Включён
      \texttt{bvhMergeAllRelems:i = 64}: треугольники модели
      собираются, сортируются по Morton-коду центроида и
      разрезаются на не менее чем 64 пространственно-плотных
      кластеров.
\end{itemize}

\paragraph{Качественная валидация.} Дополнительно для оценки
качества пространственной локализации захватываются кадры в
Nsight~Graphics с раскраской BLAS-инстансов в режиме
\texttt{Color By: Instance}. На baseline-конфигурации ожидается
«шумовая» картина из перекрывающихся прямоугольников; на
оптимизированной --- крупные пространственно-связные области,
соответствующие отдельным BLAS-кластерам.

\subsection{Сводная схема предлагаемого решения}

Итоговое отображение $\pi$ задаётся следующей схемой:

\begin{equation*}
\bigl\{M_k\bigr\}_{k=1}^{R}
\;\xrightarrow[\text{Morton-сортировка}]{\text{объединение} + }\;
\bigl\{M'_j\bigr\}_{j=1}^{K}
\;\xrightarrow[\text{add\_mesh + add\_instance}]{\text{bvh::}}\;
\bigl\{T_j\bigr\}_{j=1}^{K},
\end{equation*}
\noindent
где $\{M_k\}_{k=1}^{R}$ --- исходные наборы треугольников модели
(материальные \texttt{RElem}-ы), $\{M'_j\}_{j=1}^{K}$ --
пространственно-компактные кластеры после объединения и
Morton-сортировки, а $\{T_j\}_{j=1}^{K}$ --- построенные драйвером
BLAS этих кластеров. Качество результата оценивается по медианному
GPU-времени \texttt{rtsm::do\_trace} в \texttt{testGI}
и качественной картиной Acceleration~Structure в Nsight~Graphics;
конкретные артефакты приведены в экспериментальной части.

\subsection{Параметры алгоритма и выбор значений}

Метод имеет один пользовательский параметр:
\textbf{$K_{\text{req}}$} (\texttt{bvhMergeAllRelems})~--- желаемое число
кластеров на модель. Для целевой модели \texttt{House}
($\sim 1{,}2$~млн треугольников) выбрано $K_{\text{req}}\!=\!64$ как
оптимум по результатам sweep в экспериментальной части.
Эффективное $K$ автоматически поднимается до
$\lceil N/T_{\max}\rceil\!=\!56$ при меньших запросах;
при $K_{\text{req}}\!\ge\!64$ используется запрошенное значение.

\subsection{Выводы по разделу}

Таким образом, задача декомпозирована на три подзадачи. Анализ
pipeline формирования BLAS выявил единственную доступную без модификации драйвера
точку вмешательства --- стадию
\texttt{RElem[]} и \texttt{RaytraceGeometryDescription[]}; на её
основе построен алгоритм глобальной пространственной кластеризации
со сложностью $O(N\log N)$ и сохранением множества треугольников по
построению; определён способ валидации --- прямой замер
\texttt{rtsm::do\_trace} в \texttt{testGI} и проверка локализации BLAS в
Nsight~Graphics. Метод и его единственный параметр $K_{\text{req}}$
готовы к экспериментальной проверке в практической части.
